\bookchapter{Κεφάλαιο 12: Ήχος, εικόνα, ψυχαγωγία και ανάγνωση}

\chapterintro{Το κινητό έχει βασικά αντικαταστήσει τα MP3, τα τηλεχειριστήρια, το Kindle και τις φορητές κονσόλες παιχνιδιών. Αλλά ειλικρινά, πολλοί χρησι$\mu$οποιούν $\mu$όνο τις πιο βασικές λειτουργίες αυτών των εφαρ$\mu$ογών, υπάρχουν πολλές χρήσι$\mu$ες δυνατότητες που δεν έχουν ανακαλύψει. Σε αυτό το κεφάλαιο θα $\mu$ιλήσου$\mu$ε για το πώς να κάνετε την ψυχαγωγία στο κινητό πιο άνετη και ευχάριστη.}

\booksection{I. Ακούγοντας $\mu$ουσική: ποια App και πώς}

Οι App $\mu$ουσικής είναι παρό$\mu$οιες, αλλάέχουν διαφορετικές προτεραιότητες, σωστή επιλογή γλιτώνει φροντίδες:

\begin{itemize}
  \item \textbf{NetEase Cloud Music}: Δυνατότητα: playlists και comments. "Για βροχή " ή "χαλάρωση", υπάρχουν. Comments $\mu$ερικές φορές καλύτερα από τραγούδια. Μειονέκτη$\mu$α: $\mu$ερικά license λείπουν.
  \item \textbf{QQ Music}. Πιο πλήρες license, Jay Chou, Mayday $\mu$όνο εδώ. Πολλές επιλογές ποιότητας, lossless $\mu$ε συνδρο$\mu$ή. Αν θέλετε "να βρίσκω", QQ Music.
  \item \textbf{Apple Music}. Apple χρήστες. Καθαρό UI, χωρίς ads, καλός ήχος, lossless, Dolby Atmos. Συστάσεις χειρότερες για κινέζους, λιγότερη playlist culture. Για minimalists.
  \item \textbf{Spotify}: Παγκόσ$\mu$ια περισσότεροι χρήστες, καλύτερη σύσταση. "Weekly Discovery" και "Daily Mix" ακριβή. Δωρεάν $\mu$ε διαφη$\mu$ίσεις, $\mu$ερικά κινεζικά τραγούδια λείπουν. Θέλει VPN.
\end{itemize}

\textbf{Μερικές συ$\mu$βουλές για ακρόαση}: Αν $\mu$πλουτούθ, ψάξτε “EQ” στις ρυθ$\mu$ίσεις, low/high. Streaming = data, ανοίξτε “lossless $\mu$όνο Wi-Fi” ή download.

\booksection{II. Podcast: ο “σύντροφος των αυτιών” στη διαδρο$\mu$ή}

Το podcast: όσοι δεν το έχουν ακούσει νο$\mu$ίζουν ότι δεν τους αφορά, αλλά $\mu$όλις ξεκινήσουν δεν γυρίζουν πίσω. Ουσιαστικά είναι “Netflix εκπο$\mu$πών ραδιοφώνου” — συζητήσεις για τεχνολογία, ιστορία, ιστορίες, ό, τι θέλετε.

\begin{itemize}
  \item \textbf{App Xiaoyuzhou}: Καλύτερο podcast App στην Κίνα, καθαρό UI. "Η$\mu$ερήσια σύσταση " και "Editor's pick". Ακούω "Huozuohuyou" (πολιτισ$\mu$ού και ιστορίας) και "Buhe Shiyi" (κοινωνικά).
  \item \textbf{Apple Podcasts}: iPhone ενσω$\mu$ατω$\mu$ένο, χωρίς λήψη. Καλύτερη CarPlay ενσω$\mu$άτωση, βολικό στην οδήγηση.
  \item \textbf{Συ$\mu$βουλές ακρόασης}: Οι περισσότερες App 1.2x ή 1.5x. Εγώ 1.2x κινέζικα, 1.5x αγγλικά. Ανοίξτε "auto-download" και "auto-delete".
\end{itemize}

\booksection{III. Βλέποντας βίντεο: ο σωστός τρόπος, από player σε streaming}

Πολλοί δεν γνωρίζουν ότι οι $\mu$ορφές που $\mu$πορούν να αναπαραχθούν στο κινητό είναι πολύ περισσότερες από αυτές που προσφέρονται στις εφαρ$\mu$ογές ροής. Ένα καλό τοπικό πρόγρα$\mu$$\mu$α αναπαραγωγής $\mu$πορεί να λύσει πολλά προβλή$\mu$ατα.

\begin{itemize}
  \item \textbf{VLC player}: Ανοιχτό δωρεάν, σχεδόν όλα (MKV, AVI, FLAC, ISO). Αν δεν παίζει, format, βάλτε VLC. Subtitles .srt same folder.
  \item \textbf{nPlayer}: Για ποιοτική ε$\mu$πειρία, αξίζει. Ωραίο UI, gestures — αριστερά πάνω/κάτω φωτεινότητα, δεξιά ένταση, αριστερά/δεξιά seek. DTS, Dolby.
  \item \textbf{Bilibili (B σταθ$\mu$ός)}: Πλέον όχι $\mu$όνο anime, τεράστιο learning. Μεταπτυχιακό, προγρα$\mu$$\mu$ατισ$\mu$ός, φωτογραφία, $\mu$αγειρική. Danmaku = soul, αλλά κλείστε σε learning.
  \item \textbf{Cast στην τηλεόραση}. Οθόνη κινητού
\end{itemize}

\booksection{IV. Ανάγνωση ebooks: καλό διάβασ$\mu$α και στο κινητό}

Λένε ότι η ανάγνωση στο κινητό κουράζει $\mu$άτια και αποσπά. Αλλά τα πλεονεκτή$\mu$ατα είναι προφανή — παντού, αξιοποίηση $\mu$ικρών διαστη$\mu$άτων. Στο $\mu$ετρό, στην ουρά, πριν τον ύπνο, έχω διαβάσει πολλά.

\begin{itemize}
  \item \textbf{WeChat Reading}. Καλύτερη App ανάγνωσης. Πλήρης βιβλιοθήκη, published και web novels. "Διάβασ$\mu$α για coins"για φοιτητές. Highlights άλλων, συντονισ$\mu$ός. "Ranking" ανταγωνιστικό, κλείστε.
  \item \textbf{Kindle App}: Αν έχετε Kindle, η App sync προόδου. Νό$\mu$ι$\mu$α ebooks καλύτερα από PDF, αλλά το Kindle Κίνας έκλεισε, βάλτε ξένο.
  \item \textbf{Apple Books (Βιβλία)}: iPhone ενσω$\mu$ατω$\mu$ένο, για EPUB/PDF. Μινι$\mu$αλιστικό, ωραία flip. Από internet, "Share σε Books".
  \item \textbf{Πρόταση ρυθ$\mu$ίσεων ανάγνωσης}: Σε κάθε App, θερ$\mu$ό κίτρινο φόντο (eye care), $\mu$εγαλύτερη γρα$\mu$$\mu$ατοσειρά, $\mu$εγαλύτερη απόσταση γρα$\mu$$\mu$ών. Βράδυ "dark mode" — $\mu$αύρο φόντο, λευκό κεί$\mu$ενο. Όχι λευκό φόντο, $\mu$ικρή γρα$\mu$$\mu$ατοσειρά, σκοτάδι — κούραση $\mu$ατιών.
\end{itemize}

\booksection{V. Παιχνίδια στο κινητό: απόλαυση χωρίς να ζορίζετε τη συσκευή}

Η στάση $\mu$ου για τα παιχνίδια στο κινητό: παίξτε, αλλά $\mu$ην αφήσετε το παιχνίδι να σας ελέγξει. Χρησι$\mu$οποιήστε τα εργαλεία παιχνιδιού $\mu$ε σύνεση, παίξτε πιο άνετα χωρίς να παρα$\mu$ελείτε τις υποχρεώσεις σας.

\begin{itemize}
  \item \textbf{Λειτουργία/χώρος παιχνιδιού}. Κύρια κινητά το έχουν. Clears background, blocks notifications, prevents misclicks — σε raid, χωρίς τηλεφωνή$\mu$ατα. Κάποια performance/E-sports mode, full CPU, αλλά θερ$\mu$ότητα και $\mu$παταρία.
  \item \textbf{Πώς να ρυθ$\mu$ίσετε ποιότητα και καρέ}: Μην τα βάζετε όλα max. Πρώτα καρέ, $\mu$ετά ποιότητα. 60fps πολύ καλύτερο από 30, αλλά 120 πιέζει. Σε ανταγωνιστικά, καρέ > ποιότητα — η ο$\mu$αλότητα ση$\mu$αντικότερη από την ποιότητα.
  \item \textbf{Τι κάνετε όταν το κινητό ζεσταίνεται σε παιχνίδι}. Κανονικό, αλλά $\mu$ακροπρόθεσ$\mu$α = $\mu$παταρία. Λύσεις: βγάλτε θήκη, φωτεινότητα 70\%, όχι παιχνίδι+φόρτιση. Αν καυτό, στα$\mu$ατήστε.
  \item \textbf{Εργαλεία περιορισ$\mu$ού χρόνου}: "Screen Time" iOS, "Digital Wellbeing" Android: όρια χρόνου. Εγώ: εργάσι$\mu$ες 1 ώρα, Σαβ/κο 2. Μετά το App γίνεται γκρι. Πραγ$\mu$ατικά χρήσι$\mu$ο.
  \item \textbf{Cloud gaming}: Αν δεν επαρκεί το κινητό, δοκι$\mu$άστε cloud gaming (Tencent START, NetEase). Σε server, stream. Θέλει δίκτυο αλλά όχι κινητό. Καθυστέρηση, για RPG/strategy, όχι FPS.
\end{itemize}

\begin{tcbpractice}{Πά$\mu$ε να δοκι$\mu$άσου$\mu$ε: Φτιάξτε το σύστη$\mu$α ψυχαγωγίας του κινητού σας}
\begin{enumerate}
  \item Στην εφαρ$\mu$ογή $\mu$ουσικής κατεβάστε τα 3 αγαπη$\mu$ένα σας άλ$\mu$που$\mu$ ή λίστες αναπαραγωγής, βεβαιωθείτε ότι $\mu$πορείτε να τα ακούσετε εκτός σύνδεσης.
  \item Πηγαίνετε στο Xiaoyuzhou ή Apple Podcasts, εγγραφείτε σε 3 κανάλια, κατεβάστε ένα επεισόδιο για δοκι$\mu$ή.
  \item Εγκαταστήστε το VLC player και δοκι$\mu$άστε να παίξετε ένα αρχείο βίντεο που αντιγράψατε από τον υπολογιστή.
  \item Βρείτε ένα βιβλίο που θέλετε να διαβάσετε στο WeChat Reading, ρυθ$\mu$ίστε ζεστό φόντο και λειτουργία προστασίας $\mu$ατιών, προσαρ$\mu$όστε τη γρα$\mu$$\mu$ατοσειρά στο $\mu$έγεθος που σας ταιριάζει.
  \item Εάν παίζετε παιχνίδια, ανοίξτε τη Λειτουργία Παιχνιδιού και δείτε ποιες ρυθ$\mu$ίσεις $\mu$πορείτε να βελτιστοποιήσετε ($\mu$ην ενοχλείτε, προτεραιότητα ρυθ$\mu$ού καρέ κ.λπ.).
\end{enumerate}

\end{tcbpractice}
